%
% job-abort.spec.tex
%
% Ingest job lifecycle under a shutdown abort: a Z specification of how a
% queued job's status moves between the worker, the producer, the terminals,
% the commit checkpoint and the abort, and of the tally the abort keeps
% (quarry-70zm).
% Companion to docs/design/job-abort.md, which cites the sections below by
% number, maps each schema to the code, and records the tool runs.
%
% Check:        z-spec check docs/design/job-abort.spec.tex
% Model-check:  probcli with -p MAX_INITIALISATIONS 100 -nodead (the command
%               lines and the state counts are in docs/design/job-abort.md)
%
% Must be .tex, not .md: the z-spec tools glob for .tex.
%

\documentclass[a4paper,10pt,fleqn]{article}
\usepackage[margin=1in]{geometry}
\usepackage{fuzz}
\usepackage{booktabs}

\newcommand{\field}[1]{\texttt{#1}}

\begin{document}

\title{Ingest Job Lifecycle Under Shutdown Abort: A Z Specification}
\author{Guido R, for Punt Labs}
\date{October 2026}
\maketitle


% ===================================================================
\section{Scope and modelling choices}
\label{sec:scope}
% ===================================================================

The model has \emph{two jobs} (\field{ja}, \field{jb}) on \emph{one collection
worker}, which holds at most one job at a time and runs it on a store thread.
A job is one of three shapes, fixed when it is admitted. A \emph{backfill} job
(\field{ck} true) reports a commit checkpoint, and is run again from its
transcript, so it needs no spool. A \emph{rerunnable} job (\field{ck} false,
\field{skrerun}) is the same without the checkpoint. A \emph{spooled} job
(\field{skspool}: a \field{remember} or a URL ingest) has no durable client
copy, so the abort writes it a spool record.

Five actors write a job's status or its tally:

\begin{itemize}
  \item the \emph{worker}, which starts a job and, through its terminal
    (\field{TaskTerminal}, \field{RefusalAwareTerminal}), records how the body
    ended;
  \item the \emph{producer}, which gives up on a job by withdrawing it
    (\field{JobLifecycle.withdraw}), from another thread;
  \item the \emph{store thread}, which reaches the commit checkpoint
    (\field{JobLifecycle.commit}) and then stores or fails, and which no cancel
    stops;
  \item the \emph{shutdown abort}: \field{AbortStart} (cancel the worker, take
    the queued jobs), then per job \field{Begin}, a decision, a spool thread, and
    either \field{Record} (within the wait), \field{Deadline} (an in-flight job's
    spool outlasts \field{SPOOL\_WAIT\_S}) or, later, \field{Late};
  \item the \emph{observer} of a store that finishes after the cancel, which is
    the one actor the old build lacks.
\end{itemize}

\subsection{Statuses, the lock, and the unit of atomicity}

All seven statuses a job can hold appear: \field{squeued}, \field{srunning},
\field{scommitting}, \field{scompleted}, \field{sfailed}, \field{swithdrawn},
\field{saborted}. (\field{new}, a job not yet queued, is not modelled: a
withdraw there behaves as at \field{squeued}.) Every status write in the fixed
build is one operation, and an operation is one critical section under
\field{state.claim}: no other operation interleaves with it. That is the
single-writer rule (I6), and it is the modelling unit, not a variable. A write
that the code makes \emph{without} the claim is an operation that latches
\field{wClaim}; the old build has three, and the seeded faults add a fourth.
What the model cannot show is a check-then-write race inside such an unclaimed
write, because the model makes each write atomic whether or not the code does.

\subsection{Two builds, nine seeded faults, two switches}

\field{fixed}~$=1$ is the lifecycle this design requires. \field{fixed}~$=0$ is
the code at the commit the design starts from (eaa5118): the terminals write
without the claim and over any status, the abort's kept set lacks
\field{aborted}, the abort ignores what \field{fail\_aborted} answered and
spools and counts a kept job anyway, a store that lands after the cancel is never
observed, and a spool future that is cancelled while its thread runs is reported
as a failed write; a cancel of the abort's gather escapes the queued path's
record. \field{abl}$\,\in 1\upto9$ seeds one fault into the fixed build: the
observer moves a term and not the counter (1, I2), the late correction moves a
term and not the reason (2, I5), the observer writes without the claim (3, I6), the
abort's kept set lacks \field{aborted} (4, I7), the observer leaves the spool file
(5, I11), a failed committing store leaves the term \emph{unconfirmed} (6, I9), and
the observer leaves the failure reason on a completed job (7, I10), and the loop
closes with no drain and no report (8, I12), and a settled job is counted \emph{kept}
(9, I9). The old build
never breaks I2 or I5, so faults 1 and 2 prove those witnesses can fire. The same witnesses are
checked against every build: the fixed build, with \field{abl}$=0$, must reach
none.

\subsection{Witnesses, not invariants}

A schema invariant that an operation can violate makes that operation
unsatisfiable, so ProB would report nothing rather than a violation. Each
property is a latched witness in the set \field{wit}, computed by every operation
from the state before and after it, and checked by asking ProB whether any state
with the witness in \field{wit} is reachable.

\subsection{Properties}

A job is \emph{taken} once \field{Begin} has run (\field{ab} is past
\field{atake}): between \field{AbortStart}, or the worker's handler, and
\field{Begin} the job is in no term, and the counter of \field{tnone} holds it.
A taken job is \emph{quiet} when its outcome is recorded (\field{arec}) and its
store thread has ended or never ran.

\begin{center}
\begin{tabular}{@{}cll@{}}
\toprule
 & Witness & Meaning \\
\midrule
I1 & \field{wTerm} & a quiet job has no terminal status that matches reality:
  \field{completed} without rows, \field{failed} with rows, or still queued, running or
  committing \\
I2 & \field{wCount} & a term of the tally differs from the number of jobs in
  it (so the terms do not sum to the jobs taken) \\
I3 & \field{wKept} & a withdrawn, aborted or stored job is counted in a term
  other than \emph{kept}: spooled, lost, unfinished, or needing no spool \\
I4 & \field{wLost} & a job's rows are stored and its status is \field{failed},
  unless the loop closed first and reported it \\
I5 & \field{wLate} & a failed job's reason and its term disagree (a late
  correction moved one and not the other) \\
I6 & \field{wClaim} & a status was written without the claim \\
I7 & \field{wOver} & a withdrawn, aborted or completed status was overwritten,
  or a committing job was failed while its store thread still runs \\
I8 & \field{wDisk} & a failed job's reason says what reached disk wrongly:
  spooled with no file, lost though the write may yet land or has, or no spool
  needed though a file exists \\
I9 & \field{wMatch} & a term does not match the status: \emph{unconfirmed}
  without \field{committing}; \emph{pending}, \emph{unfinished}, \emph{spooled},
  \emph{not written}, \emph{unbuilt} or \emph{needs no spool} without
  \field{failed}; \emph{kept} without withdrawn, aborted or completed \\
I10 & \field{wCompl} & a \field{completed} job still carries a failure reason \\
I11 & \field{wGhost} & a \field{completed} job has a spool file on disk, unless the
  loop closed first and reported it \\
I12 & \field{wSilent} & the loop has closed while a job's store, or a spool write
  for a \field{completed} job, was still running, and the close did not report it \\
\bottomrule
\end{tabular}
\end{center}

I2 is also nearly structural: \field{tl} is a function from jobs to terms, so
every job is in exactly one term. The counters \field{cnt} are what the code
keeps, and \field{wCount} checks them against \field{tl}.

\subsection{The loop closes}

The observer is a callback on the event loop, and the loop does not outlive every
thread. \field{IngestQueue.aclose} can return while a detached store is still
running, and \field{asyncio.run} then closes the loop and cancels the observer: the
store's end is never seen, the job stays \field{failed} with its rows, its term is
never corrected, and a spool file is never removed. \field{LoopClose} models the
close. The fixed build drains before it, waiting for the detached stores and the
pending spool removals until the shutdown budget expires (\field{BudgetExpires},
which may fire at any time after the abort starts), and then \emph{reports} every
job whose store or spool write is still running (\field{rp}): a warning names its
task, and the tally that is logged is the corrected one. After \field{LoopClose}
no observer, removal or late callback runs, and a store that ends is
\field{StoreEndsUnobserved}: the status is left as it was. I4 and I11 therefore
hold \emph{for stores that end before the close}; for the others the invariant is
I12, that none is silent. A store that never ends, and the budget's length, are not
modelled.

A late store after a spool record was written would leave the record on disk: a
\emph{ghost}. The fixed build forbids it (I11). When the observer completes a
job it removes the job's spool file, and a spool thread that ends after the job
has stored removes the file it wrote (\field{SpoolEndsRemoved}). The model makes
each removal atomic with the step that causes it; the code cannot, so a crash
between the store and the removal can still leave a file, which the daemon's
start-up sweep does not touch (it removes temps, not records). Whether replaying
such a record is harmless is not modelled, and no replay code exists to check it.

\subsection{Assumption: the store has its own executor}

The store thread and the spool thread never wait for each other. The code must
make that true: the spool writes run on the shared \field{DaemonThreads} pool of
four threads, so stores must run on a separate executor, or four busy
collections would starve every spool write and the abort's wait would expire on
all of them. The model has no thread pool; this is the assumption that lets it
treat the two kinds of thread as independent.

\subsection{Non-goals}

\begin{itemize}
  \item Time. \field{Deadline} fires whenever its guard holds; that it fires
    after \field{SPOOL\_WAIT\_S} is not modelled. So the budget is a claim about
    structure (an in-flight job's term can leave \emph{pending}), not about
    seconds.
  \item A queued job's spool has no deadline, though a cancel of the abort's
    gather does end its wait (\field{QueueCancel}): its term stays \emph{pending}
    until the spool ends, which the summary reports as not yet accounted for.
    Whether to bound it is a design question, not a model result.
  \item More than two jobs, more than one collection worker, and the embed gate
    shared between workers (the gate is one place to wait, \field{gate}).
  \item The order of the queue, and the producer's wait loop; the producer is the
    \field{Withdraw} operations alone.
  \item A second cancel during the in-flight wait: it only repeats the cancel,
    which a spool-wait loop absorbs; it changes no state the model keeps.
  \item The spool record's bytes, its recoverability flag (recoverable or
    not), and the spool directory sweep. Both spooled outcomes are one term.
  \item The summary's text, and the done-callback's logging.
  \item A store thread that never ends. It is a state with the thread not quiet,
    and a job in it stays \emph{unconfirmed} or \emph{pending} by design; the
    liveness check below is \emph{possibility}, not inevitability.
\end{itemize}

% ===================================================================
\section{Types and constants}
\label{sec:types}
% ===================================================================

\begin{zed}
ZBOOL ::= ztrue | zfalse
\end{zed}

\begin{zed}
JOB ::= ja | jb
\end{zed}

The seven statuses of a job.

\begin{zed}
STATUS ::= squeued | srunning | scommitting | scompleted | sfailed
         | swithdrawn | saborted
\end{zed}

Where the worker has the job: waiting in the collection's queue
(\field{inq}); dequeued, waiting for the embed gate (\field{gate}); in its
terminal, awaiting the store thread (\field{body}); taken out of the queue by the
abort (\field{qtaken}); or no longer the worker's, because it finished, was
skipped, or was cancelled (\field{out}). A job that is \field{out} while its store
thread runs is \emph{detached}.

\begin{zed}
LOC ::= inq | gate | body | qtaken | out
\end{zed}

The store thread. \field{t0}: not started. \field{tpre}: running, before the
commit checkpoint. \field{tcommit}: past the checkpoint (a backfill job only).
\field{tstored}: the rows are in the table. \field{tstop}: ended with none (it
stopped at the checkpoint, was refused, or failed after the delete).

\begin{zed}
TS ::= t0 | tpre | tcommit | tstored | tstop
\end{zed}

The abort's progress on a job. \field{an}: not taken. \field{atake}: taken out
of the queue, or the worker's handler has it, and \field{Begin} has not run.
\field{abegun}: failed provisionally and counted pending. \field{aspool}: its spool
write has started. \field{arec}: its outcome is recorded.

\begin{zed}
AB ::= an | atake | abegun | aspool | arec
\end{zed}

The spool thread.

\begin{zed}
SP ::= s0 | srun | sok | sfail | sunb
\end{zed}

\begin{zed}
SK ::= skrerun | skspool
\end{zed}

The tally's terms. \field{tnone}: not taken (the counter holds the jobs not yet
counted, so the terms sum to two). \field{tpend}: taken, outcome not yet known.
\field{tspooled}, \field{tnotw} (the write failed), \field{tunb} (the record
could not be built), \field{tunfin} (the in-flight wait ended with the write still
running), \field{tnospool} (it needs none). \field{tkept}: the abort left the job
as it is: withdrawn, aborted by its producer, or stored. \field{tunconf}: a
committing job, whose store has not ended.

\begin{zed}
TERM ::= tnone | tpend | tspooled | tnospool | tnotw | tunb | tunfin
       | tkept | tunconf
\end{zed}

The reason a job's status carries. \field{rterm}: the terminal's own (cancelled,
refused, failed). \field{rprov}: the abort's provisional ``spool unfinished at
shutdown'', which is also the final reason of an unfinished job.

\begin{zed}
RS ::= rnone | rterm | rprov | rspooled | rnospool | rnotw | runb
\end{zed}

\begin{zed}
PHASE ::= pre | started
\end{zed}

\begin{zed}
WIT ::= wTerm | wCount | wKept | wLost | wLate | wClaim | wOver | wDisk
      | wMatch | wCompl | wGhost | wSilent
\end{zed}

\begin{axdef}
fixed, abl : \nat
\where
fixed \in 0 \upto 1 \\
abl \in 0 \upto 9 \\
abl > 0 \implies fixed = 1 \\
fixed = 1 \\
abl = 0
\end{axdef}

\begin{axdef}
ck : JOB \fun ZBOOL \\
sk : JOB \fun SK
\where
\forall j : JOB @ ck~j = ztrue \implies sk~j = skrerun
\end{axdef}

The reason that goes with each term, and the term and reason a spool outcome
gives. A kept or unconfirmed job carries no abort reason.

\begin{axdef}
ReasonOf : TERM \pfun RS \\
TermOf : SP \pfun TERM \\
RsOf : SP \pfun RS
\where
ReasonOf = \{ tpend \mapsto rprov, tunfin \mapsto rprov,
  tspooled \mapsto rspooled,
  tnotw \mapsto rnotw, tunb \mapsto runb \} \\
TermOf = \{ sok \mapsto tspooled, sfail \mapsto tnotw, sunb \mapsto tunb \} \\
RsOf = \{ sok \mapsto rspooled, sfail \mapsto rnotw, sunb \mapsto runb \}
\end{axdef}


% ===================================================================
\section{State}
\label{sec:state}
% ===================================================================

\begin{schema}{Run}
st : JOB \fun STATUS \\
loc : JOB \fun LOC \\
ts : JOB \fun TS \\
ab : JOB \fun AB \\
sp : JOB \fun SP \\
tl : JOB \fun TERM \\
rs : JOB \fun RS \\
fl, fc, rp : JOB \fun ZBOOL \\
cnt : TERM \fun (0 \upto 2) \\
ph : PHASE \\
lc, bx : ZBOOL \\
wit : \power WIT
\end{schema}

\begin{itemize}
  \item \field{st}, \field{rs}: the job's status and the reason its \field{error}
    carries.
  \item \field{loc}, \field{ts}: the worker's hold on it and its store thread.
  \item \field{ab}, \field{sp}: the abort's progress on it and its spool thread.
    The spool file is on disk exactly when \field{sp}~$=$~\field{sok}.
  \item \field{fl}: the job was in flight (the worker held it) when it was
    taken. \field{fc}: its spool future was cancelled while the thread ran.
  \item \field{tl}, \field{cnt}: the job's term and the counters the code keeps;
    the code writes them together or the model reports it.
  \item \field{ph}: the abort has started. \field{lc}: the event loop has
    closed, so no done-callback, observer or removal runs any more. \field{bx}:
    the shutdown budget for the drain before the close has expired.
    \field{rp}: the close reported the job (a warning naming its task): its store
    was still running, or a spool write for a completed job was.
\end{itemize}

\begin{schema}{Init}
Run'
\where
\forall j : JOB @ st'~j = squeued \land loc'~j = inq \land ts'~j = t0 \land
  ab'~j = an \land sp'~j = s0 \land tl'~j = tnone \land rs'~j = rnone \land
  fl'~j = zfalse \land fc'~j = zfalse \land rp'~j = zfalse \\
cnt' = \{ tnone \mapsto 2, tpend \mapsto 0, tspooled \mapsto 0, tnospool \mapsto 0,
  tnotw \mapsto 0, tunb \mapsto 0, tunfin \mapsto 0, tkept \mapsto 0,
  tunconf \mapsto 0 \} \\
ph' = pre \\
lc' = zfalse \\
bx' = zfalse \\
wit' = \emptyset
\end{schema}

Both jobs start queued. \field{ck} and \field{sk} are free.


% ===================================================================
\section{Latching the witnesses}
\label{sec:latch}
% ===================================================================

Every operation below is its own core schema conjoined with \field{Latch}, which computes \field{wit}$'$ from the state before and after the step and from nothing else, so that no operation can forget to report a violation. An operation that the code performs without holding \field{state.claim} is conjoined with \field{LatchClaim}, which adds \field{wClaim}. Both are conjoined, never invoked on their own: the operations are the schemas defined with $\defs$.

\begin{schema}{Latch}
\Delta Run
\where
wit' = wit \cup \{ w : WIT |
  (w = wTerm \land (\exists j : JOB @ ab'~j = arec \land rp'~j = zfalse \land
     ts'~j \in \{ t0, tstored, tstop \} \land
     \lnot ((st'~j = scompleted \land ts'~j = tstored) \lor
            (st'~j = sfailed \land ts'~j \neq tstored) \lor
            st'~j \in \{ swithdrawn, saborted \}))) \lor
  (w = wCount \land (\exists t : TERM @
     (\# \{ j : JOB | tl'~j = t \}) \neq cnt'~t)) \lor
  (w = wKept \land (\exists j : JOB @ st'~j \in \{ swithdrawn, saborted, scompleted \} \land
     tl'~j \notin \{ tnone, tkept \})) \lor
  (w = wLost \land (\exists j : JOB @ ts'~j = tstored \land st'~j = sfailed \land
     rp'~j = zfalse)) \lor
  (w = wLate \land (\exists j : JOB @ st'~j = sfailed \land
     ((tl'~j \in \dom ReasonOf \land rs'~j \neq ReasonOf~(tl'~j)) \lor
      (tl'~j = tnospool \land rs'~j \notin \{ rnospool, rterm \})))) \lor
  (w = wDisk \land (\exists j : JOB @ st'~j = sfailed \land
     ((rs'~j = rspooled \land sp'~j \neq sok) \lor
      (rs'~j \in \{ rnotw, runb \} \land sp'~j \in \{ srun, sok \}) \lor
      (rs'~j = rnospool \land sp'~j = sok)))) \lor
  (w = wOver \land (\exists j : JOB @
     (st~j \in \{ swithdrawn, saborted, scompleted \} \land st'~j \neq st~j) \lor
     (st~j = scommitting \land st'~j = sfailed \land ts'~j = tcommit))) \lor
  (w = wMatch \land (\exists j : JOB @
     (tl'~j = tunconf \land st'~j \neq scommitting) \lor
     (tl'~j \in \{ tpend, tunfin, tspooled, tnotw, tunb, tnospool \} \land
        st'~j \neq sfailed) \lor
     (tl'~j = tkept \land st'~j \notin \{ swithdrawn, saborted, scompleted \}))) \lor
  (w = wCompl \land (\exists j : JOB @ st'~j = scompleted \land rs'~j \neq rnone)) \lor
  (w = wGhost \land (\exists j : JOB @ st'~j = scompleted \land sp'~j = sok \land
     rp'~j = zfalse)) \lor
  (w = wSilent \land lc' = ztrue \land (\exists j : JOB @ rp'~j = zfalse \land
     (ts'~j \in \{ tpre, tcommit \} \lor (st'~j = scompleted \land sp'~j = srun)))) \}
\end{schema}

\begin{schema}{LatchClaim}
\Delta Run
\where
wit' = wit \cup \{ w : WIT |
  (w = wTerm \land (\exists j : JOB @ ab'~j = arec \land rp'~j = zfalse \land
     ts'~j \in \{ t0, tstored, tstop \} \land
     \lnot ((st'~j = scompleted \land ts'~j = tstored) \lor
            (st'~j = sfailed \land ts'~j \neq tstored) \lor
            st'~j \in \{ swithdrawn, saborted \}))) \lor
  (w = wCount \land (\exists t : TERM @
     (\# \{ j : JOB | tl'~j = t \}) \neq cnt'~t)) \lor
  (w = wKept \land (\exists j : JOB @ st'~j \in \{ swithdrawn, saborted, scompleted \} \land
     tl'~j \notin \{ tnone, tkept \})) \lor
  (w = wLost \land (\exists j : JOB @ ts'~j = tstored \land st'~j = sfailed \land
     rp'~j = zfalse)) \lor
  (w = wLate \land (\exists j : JOB @ st'~j = sfailed \land
     ((tl'~j \in \dom ReasonOf \land rs'~j \neq ReasonOf~(tl'~j)) \lor
      (tl'~j = tnospool \land rs'~j \notin \{ rnospool, rterm \})))) \lor
  (w = wDisk \land (\exists j : JOB @ st'~j = sfailed \land
     ((rs'~j = rspooled \land sp'~j \neq sok) \lor
      (rs'~j \in \{ rnotw, runb \} \land sp'~j \in \{ srun, sok \}) \lor
      (rs'~j = rnospool \land sp'~j = sok)))) \lor
  (w = wOver \land (\exists j : JOB @
     (st~j \in \{ swithdrawn, saborted, scompleted \} \land st'~j \neq st~j) \lor
     (st~j = scommitting \land st'~j = sfailed \land ts'~j = tcommit))) \lor
  (w = wMatch \land (\exists j : JOB @
     (tl'~j = tunconf \land st'~j \neq scommitting) \lor
     (tl'~j \in \{ tpend, tunfin, tspooled, tnotw, tunb, tnospool \} \land
        st'~j \neq sfailed) \lor
     (tl'~j = tkept \land st'~j \notin \{ swithdrawn, saborted, scompleted \}))) \lor
  (w = wCompl \land (\exists j : JOB @ st'~j = scompleted \land rs'~j \neq rnone)) \lor
  (w = wGhost \land (\exists j : JOB @ st'~j = scompleted \land sp'~j = sok \land
     rp'~j = zfalse)) \lor
  (w = wSilent \land lc' = ztrue \land (\exists j : JOB @ rp'~j = zfalse \land
     (ts'~j \in \{ tpre, tcommit \} \lor (st'~j = scompleted \land sp'~j = srun)))) \lor w = wClaim \}
\end{schema}

% ===================================================================
\section{The producer, the worker and the store thread}
\label{sec:worker}
% ===================================================================

Operations in this section happen with or without a shutdown. The store thread
runs on its own: the checkpoint and the end of the store are its steps, and
whether the worker still awaits it (\field{loc}~$=$~\field{body}) or a cancel
detached it (\field{loc}~$=$~\field{out}) decides who records the result.
Every operation spells out what it leaves unchanged.

The producer withdraws a job that has not started (\field{JobLifecycle.withdraw},
under the claim).

\begin{schema}{WithdrawQueuedCore}
\Delta Run \\ j? : JOB
\where
st~j? = squeued \\
st' = st \oplus \{ j? \mapsto swithdrawn \} \\
loc' = loc \\
ts' = ts \\
ab' = ab \\
sp' = sp \\
tl' = tl \\
rs' = rs \\
fl' = fl \\
fc' = fc \\
rp' = rp \\
cnt' = cnt \\
ph' = ph \\
lc' = lc \\
bx' = bx
\end{schema}

\begin{zed}
WithdrawQueued \defs WithdrawQueuedCore \land Latch
\end{zed}

The producer gives up on a job that has started: it is \field{aborted}, and a
job with a commit checkpoint will stop at it.

\begin{schema}{WithdrawRunningCore}
\Delta Run \\ j? : JOB
\where
st~j? = srunning \\
st' = st \oplus \{ j? \mapsto saborted \} \\
loc' = loc \\
ts' = ts \\
ab' = ab \\
sp' = sp \\
tl' = tl \\
rs' = rs \\
fl' = fl \\
fc' = fc \\
rp' = rp \\
cnt' = cnt \\
ph' = ph \\
lc' = lc \\
bx' = bx
\end{schema}

\begin{zed}
WithdrawRunning \defs WithdrawRunningCore \land Latch
\end{zed}

The worker takes the next queued job. It holds at most one.

\begin{schema}{DequeueCore}
\Delta Run \\ j? : JOB
\where
ph = pre \\
loc~j? = inq \\
\forall k : JOB @ loc~k \notin \{ gate, body \} \\
st' = st \\
loc' = loc \oplus \{ j? \mapsto gate \} \\
ts' = ts \\
ab' = ab \\
sp' = sp \\
tl' = tl \\
rs' = rs \\
fl' = fl \\
fc' = fc \\
rp' = rp \\
cnt' = cnt \\
ph' = ph \\
lc' = lc \\
bx' = bx
\end{schema}

\begin{zed}
Dequeue \defs DequeueCore \land Latch
\end{zed}

\field{JobLifecycle.start}: a queued job becomes running and its store thread
begins.

\begin{schema}{StartCore}
\Delta Run \\ j? : JOB
\where
ph = pre \\
loc~j? = gate \\
st~j? = squeued \\
st' = st \oplus \{ j? \mapsto srunning \} \\
loc' = loc \oplus \{ j? \mapsto body \} \\
ts' = ts \oplus \{ j? \mapsto tpre \} \\
ab' = ab \\
sp' = sp \\
tl' = tl \\
rs' = rs \\
fl' = fl \\
fc' = fc \\
rp' = rp \\
cnt' = cnt \\
ph' = ph \\
lc' = lc \\
bx' = bx
\end{schema}

\begin{zed}
Start \defs StartCore \land Latch
\end{zed}

A job withdrawn while it waited is never started; the worker lets it go.

\begin{schema}{SkipCore}
\Delta Run \\ j? : JOB
\where
ph = pre \\
loc~j? = gate \\
st~j? \neq squeued \\
st' = st \\
loc' = loc \oplus \{ j? \mapsto out \} \\
ts' = ts \\
ab' = ab \\
sp' = sp \\
tl' = tl \\
rs' = rs \\
fl' = fl \\
fc' = fc \\
rp' = rp \\
cnt' = cnt \\
ph' = ph \\
lc' = lc \\
bx' = bx
\end{schema}

\begin{zed}
Skip \defs SkipCore \land Latch
\end{zed}

The commit checkpoint, reported by the store thread: \field{JobLifecycle.commit}
claims it only while the job is still running.

\begin{schema}{CommitCore}
\Delta Run \\ j? : JOB
\where
ck~j? = ztrue \\
ts~j? = tpre \\
st~j? = srunning \\
st' = st \oplus \{ j? \mapsto scommitting \} \\
loc' = loc \\
ts' = ts \oplus \{ j? \mapsto tcommit \} \\
ab' = ab \\
sp' = sp \\
tl' = tl \\
rs' = rs \\
fl' = fl \\
fc' = fc \\
rp' = rp \\
cnt' = cnt \\
ph' = ph \\
lc' = lc \\
bx' = bx
\end{schema}

\begin{zed}
Commit \defs CommitCore \land Latch
\end{zed}

A job whose status is no longer running refuses the commit (the checkpoint
raises), so the thread stops before it deletes or inserts anything.

\begin{schema}{CommitRefusedCore}
\Delta Run \\ j? : JOB
\where
ck~j? = ztrue \\
ts~j? = tpre \\
st~j? \neq srunning \\
st' = st \\
loc' = loc \\
ts' = ts \oplus \{ j? \mapsto tstop \} \\
ab' = ab \\
sp' = sp \\
tl' = tl \\
rs' = rs \\
fl' = fl \\
fc' = fc \\
rp' = rp \\
cnt' = cnt \\
ph' = ph \\
lc' = lc \\
bx' = bx
\end{schema}

\begin{zed}
CommitRefused \defs CommitRefusedCore \land Latch
\end{zed}

The store thread finishes and the rows are in the table, while the worker still
awaits it (its status write is the worker's), or whatever the build.

\begin{schema}{StoreEndsCore}
\Delta Run \\ j? : JOB
\where
((ck~j? = zfalse \land ts~j? = tpre) \lor (ck~j? = ztrue \land ts~j? = tcommit)) \\
loc~j? = body \\
st' = st \\
loc' = loc \\
ts' = ts \oplus \{ j? \mapsto tstored \} \\
ab' = ab \\
sp' = sp \\
tl' = tl \\
rs' = rs \\
fl' = fl \\
fc' = fc \\
rp' = rp \\
cnt' = cnt \\
ph' = ph \\
lc' = lc \\
bx' = bx
\end{schema}

\begin{zed}
StoreEnds \defs StoreEndsCore \land Latch
\end{zed}

The store thread ends without rows (a refusal, an insert that failed after the
delete), while the worker still awaits it.

\begin{schema}{StoreFailsCore}
\Delta Run \\ j? : JOB
\where
((ck~j? = zfalse \land ts~j? = tpre) \lor (ck~j? = ztrue \land ts~j? = tcommit)) \\
loc~j? = body \\
st' = st \\
loc' = loc \\
ts' = ts \oplus \{ j? \mapsto tstop \} \\
ab' = ab \\
sp' = sp \\
tl' = tl \\
rs' = rs \\
fl' = fl \\
fc' = fc \\
rp' = rp \\
cnt' = cnt \\
ph' = ph \\
lc' = lc \\
bx' = bx
\end{schema}

\begin{zed}
StoreFails \defs StoreFailsCore \land Latch
\end{zed}


% ===================================================================
\section{The store thread after a cancel}
\label{sec:detached}
% ===================================================================

The old build never looks at a store thread again once its worker is cancelled.
The fixed build keeps the thread's future with the job and observes its end: that
observer is what makes a store that lands after the cancel report
\field{completed} (I4) without a second writer.

\emph{Old build.} The worker is gone (a cancel detached it) and nothing watches
the store thread: the rows land and the status stays what the abort left.

\begin{schema}{StoreEndsOldCore}
\Delta Run \\ j? : JOB
\where
fixed = 0 \\
((ck~j? = zfalse \land ts~j? = tpre) \lor (ck~j? = ztrue \land ts~j? = tcommit)) \\
loc~j? = out \\
st' = st \\
loc' = loc \\
ts' = ts \oplus \{ j? \mapsto tstored \} \\
ab' = ab \\
sp' = sp \\
tl' = tl \\
rs' = rs \\
fl' = fl \\
fc' = fc \\
rp' = rp \\
cnt' = cnt \\
ph' = ph \\
lc' = lc \\
bx' = bx
\end{schema}

\begin{zed}
StoreEndsOld \defs StoreEndsOldCore \land Latch
\end{zed}

\emph{Old build.} The detached thread ends without rows; nobody observes it.

\begin{schema}{StoreFailsOldCore}
\Delta Run \\ j? : JOB
\where
fixed = 0 \\
((ck~j? = zfalse \land ts~j? = tpre) \lor (ck~j? = ztrue \land ts~j? = tcommit)) \\
loc~j? = out \\
st' = st \\
loc' = loc \\
ts' = ts \oplus \{ j? \mapsto tstop \} \\
ab' = ab \\
sp' = sp \\
tl' = tl \\
rs' = rs \\
fl' = fl \\
fc' = fc \\
rp' = rp \\
cnt' = cnt \\
ph' = ph \\
lc' = lc \\
bx' = bx
\end{schema}

\begin{zed}
StoreFailsOld \defs StoreFailsOldCore \land Latch
\end{zed}

\emph{Fixed build.} The detached store thread's own completion is observed (a
done-callback on the future the job keeps): its rows are stored, so the job is
\field{completed}, its reason is cleared, and its tally term moves to
\emph{kept}, in one critical section under the claim.

\begin{schema}{StoreEndsObservedCore}
\Delta Run \\ j? : JOB
\where
fixed = 1 \\
abl \neq 1 \\
abl \neq 3 \\
abl \neq 5 \\
((ck~j? = zfalse \land ts~j? = tpre) \lor (ck~j? = ztrue \land ts~j? = tcommit)) \\
loc~j? = out \\
st~j? \in \{ sfailed, scommitting \} \\
tl~j? \in \{ tpend, tspooled, tnospool, tnotw, tunb, tunfin, tunconf \} \\
lc = zfalse \\
st' = st \oplus \{ j? \mapsto scompleted \} \\
loc' = loc \\
ts' = ts \oplus \{ j? \mapsto tstored \} \\
ab' = ab \oplus \{ j? \mapsto arec \} \\
(sp~j? = sok \implies sp' = sp \oplus \{ j? \mapsto s0 \}) \land (sp~j? \neq sok \implies sp' = sp) \\
tl' = tl \oplus \{ j? \mapsto tkept \} \\
rs' = rs \oplus \{ j? \mapsto rnone \} \\
fl' = fl \\
fc' = fc \\
rp' = rp \\
cnt' = cnt \oplus \{ tl~j? \mapsto cnt~(tl~j?) - 1, tkept \mapsto cnt~tkept + 1 \} \\
ph' = ph \\
lc' = lc \\
bx' = bx
\end{schema}

\begin{zed}
StoreEndsObserved \defs StoreEndsObservedCore \land Latch
\end{zed}

\emph{Seeded fault 1.} The observer moves the job's term but not the counter.

\begin{schema}{StoreEndsObservedNoCountCore}
\Delta Run \\ j? : JOB
\where
fixed = 1 \\
abl = 1 \\
((ck~j? = zfalse \land ts~j? = tpre) \lor (ck~j? = ztrue \land ts~j? = tcommit)) \\
loc~j? = out \\
st~j? \in \{ sfailed, scommitting \} \\
tl~j? \in \{ tpend, tspooled, tnospool, tnotw, tunb, tunfin, tunconf \} \\
lc = zfalse \\
st' = st \oplus \{ j? \mapsto scompleted \} \\
loc' = loc \\
ts' = ts \oplus \{ j? \mapsto tstored \} \\
ab' = ab \oplus \{ j? \mapsto arec \} \\
(sp~j? = sok \implies sp' = sp \oplus \{ j? \mapsto s0 \}) \land (sp~j? \neq sok \implies sp' = sp) \\
tl' = tl \oplus \{ j? \mapsto tkept \} \\
rs' = rs \oplus \{ j? \mapsto rnone \} \\
fl' = fl \\
fc' = fc \\
rp' = rp \\
cnt' = cnt \\
ph' = ph \\
lc' = lc \\
bx' = bx
\end{schema}

\begin{zed}
StoreEndsObservedNoCount \defs StoreEndsObservedNoCountCore \land Latch
\end{zed}

\emph{Seeded fault 3.} The observer writes the status without the claim.

\begin{schema}{StoreEndsObservedUnclaimedCore}
\Delta Run \\ j? : JOB
\where
fixed = 1 \\
abl = 3 \\
((ck~j? = zfalse \land ts~j? = tpre) \lor (ck~j? = ztrue \land ts~j? = tcommit)) \\
loc~j? = out \\
st~j? \in \{ sfailed, scommitting \} \\
tl~j? \in \{ tpend, tspooled, tnospool, tnotw, tunb, tunfin, tunconf \} \\
lc = zfalse \\
st' = st \oplus \{ j? \mapsto scompleted \} \\
loc' = loc \\
ts' = ts \oplus \{ j? \mapsto tstored \} \\
ab' = ab \oplus \{ j? \mapsto arec \} \\
(sp~j? = sok \implies sp' = sp \oplus \{ j? \mapsto s0 \}) \land (sp~j? \neq sok \implies sp' = sp) \\
tl' = tl \oplus \{ j? \mapsto tkept \} \\
rs' = rs \oplus \{ j? \mapsto rnone \} \\
fl' = fl \\
fc' = fc \\
rp' = rp \\
cnt' = cnt \oplus \{ tl~j? \mapsto cnt~(tl~j?) - 1, tkept \mapsto cnt~tkept + 1 \} \\
ph' = ph \\
lc' = lc \\
bx' = bx
\end{schema}

\begin{zed}
StoreEndsObservedUnclaimed \defs StoreEndsObservedUnclaimedCore \land LatchClaim
\end{zed}

\emph{Fixed build.} The store lands after the cancel and before the abort has
begun the job: the status is corrected now, its reason is cleared, and
\field{Begin} will find it kept.

\begin{schema}{StoreEndsObservedUnbegunCore}
\Delta Run \\ j? : JOB
\where
fixed = 1 \\
abl \neq 7 \\
((ck~j? = zfalse \land ts~j? = tpre) \lor (ck~j? = ztrue \land ts~j? = tcommit)) \\
loc~j? = out \\
st~j? \in \{ sfailed, scommitting \} \\
tl~j? = tnone \\
lc = zfalse \\
st' = st \oplus \{ j? \mapsto scompleted \} \\
loc' = loc \\
ts' = ts \oplus \{ j? \mapsto tstored \} \\
ab' = ab \\
sp' = sp \\
tl' = tl \\
rs' = rs \oplus \{ j? \mapsto rnone \} \\
fl' = fl \\
fc' = fc \\
rp' = rp \\
cnt' = cnt \\
ph' = ph \\
lc' = lc \\
bx' = bx
\end{schema}

\begin{zed}
StoreEndsObservedUnbegun \defs StoreEndsObservedUnbegunCore \land Latch
\end{zed}

\emph{Seeded fault 7.} The same, and the failure reason is left on the job.

\begin{schema}{StoreEndsObservedUnbegunKeepsReasonCore}
\Delta Run \\ j? : JOB
\where
abl = 7 \\
((ck~j? = zfalse \land ts~j? = tpre) \lor (ck~j? = ztrue \land ts~j? = tcommit)) \\
loc~j? = out \\
st~j? \in \{ sfailed, scommitting \} \\
tl~j? = tnone \\
lc = zfalse \\
st' = st \oplus \{ j? \mapsto scompleted \} \\
loc' = loc \\
ts' = ts \oplus \{ j? \mapsto tstored \} \\
ab' = ab \\
sp' = sp \\
tl' = tl \\
rs' = rs \\
fl' = fl \\
fc' = fc \\
rp' = rp \\
cnt' = cnt \\
ph' = ph \\
lc' = lc \\
bx' = bx
\end{schema}

\begin{zed}
StoreEndsObservedUnbegunKeepsReason \defs StoreEndsObservedUnbegunKeepsReasonCore \land Latch
\end{zed}

\emph{Seeded fault 5.} The observer leaves the job's spool file on disk.

\begin{schema}{StoreEndsObservedGhostCore}
\Delta Run \\ j? : JOB
\where
fixed = 1 \\
abl = 5 \\
((ck~j? = zfalse \land ts~j? = tpre) \lor (ck~j? = ztrue \land ts~j? = tcommit)) \\
loc~j? = out \\
st~j? \in \{ sfailed, scommitting \} \\
tl~j? \in \{ tpend, tspooled, tnospool, tnotw, tunb, tunfin, tunconf \} \\
lc = zfalse \\
st' = st \oplus \{ j? \mapsto scompleted \} \\
loc' = loc \\
ts' = ts \oplus \{ j? \mapsto tstored \} \\
ab' = ab \oplus \{ j? \mapsto arec \} \\
sp' = sp \\
tl' = tl \oplus \{ j? \mapsto tkept \} \\
rs' = rs \oplus \{ j? \mapsto rnone \} \\
fl' = fl \\
fc' = fc \\
rp' = rp \\
cnt' = cnt \oplus \{ tl~j? \mapsto cnt~(tl~j?) - 1, tkept \mapsto cnt~tkept + 1 \} \\
ph' = ph \\
lc' = lc \\
bx' = bx
\end{schema}

\begin{zed}
StoreEndsObservedGhost \defs StoreEndsObservedGhostCore \land Latch
\end{zed}

\emph{Fixed build.} The spool thread ends after the job has stored: its done
callback sees \field{completed} and removes the file it wrote (the write returns
the path), so no ghost record survives.

\begin{schema}{SpoolEndsRemovedCore}
\Delta Run \\ j? : JOB
\where
fixed = 1 \\
abl \neq 5 \\
sp~j? = srun \\
st~j? = scompleted \\
lc = zfalse \\
st' = st \\
loc' = loc \\
ts' = ts \\
ab' = ab \\
sp' = sp \oplus \{ j? \mapsto s0 \} \\
tl' = tl \\
rs' = rs \\
fl' = fl \\
fc' = fc \\
rp' = rp \\
cnt' = cnt \\
ph' = ph \\
lc' = lc \\
bx' = bx
\end{schema}

\begin{zed}
SpoolEndsRemoved \defs SpoolEndsRemovedCore \land Latch
\end{zed}

\emph{Seeded fault 6.} The same, and the term stays \emph{unconfirmed}.

\begin{schema}{StoreFailsObservedCommittingStaleCore}
\Delta Run \\ j? : JOB
\where
abl = 6 \\
fixed = 1 \\
((ck~j? = zfalse \land ts~j? = tpre) \lor (ck~j? = ztrue \land ts~j? = tcommit)) \\
loc~j? = out \\
st~j? = scommitting \\
tl~j? = tunconf \\
lc = zfalse \\
st' = st \oplus \{ j? \mapsto sfailed \} \\
loc' = loc \\
ts' = ts \oplus \{ j? \mapsto tstop \} \\
ab' = ab \\
sp' = sp \\
tl' = tl \\
rs' = rs \oplus \{ j? \mapsto rnospool \} \\
fl' = fl \\
fc' = fc \\
rp' = rp \\
cnt' = cnt \\
ph' = ph \\
lc' = lc \\
bx' = bx
\end{schema}

\begin{zed}
StoreFailsObservedCommittingStale \defs StoreFailsObservedCommittingStaleCore \land Latch
\end{zed}

\emph{Fixed build.} The detached store lands on a status the producer or an
earlier correction already settled; the status is left alone.

\begin{schema}{StoreEndsKeptCore}
\Delta Run \\ j? : JOB
\where
fixed = 1 \\
((ck~j? = zfalse \land ts~j? = tpre) \lor (ck~j? = ztrue \land ts~j? = tcommit)) \\
loc~j? = out \\
st~j? \in \{ swithdrawn, saborted, scompleted \} \\
lc = zfalse \\
st' = st \\
loc' = loc \\
ts' = ts \oplus \{ j? \mapsto tstored \} \\
ab' = ab \\
sp' = sp \\
tl' = tl \\
rs' = rs \\
fl' = fl \\
fc' = fc \\
rp' = rp \\
cnt' = cnt \\
ph' = ph \\
lc' = lc \\
bx' = bx
\end{schema}

\begin{zed}
StoreEndsKept \defs StoreEndsKeptCore \land Latch
\end{zed}

\emph{Fixed build.} The detached store of a committing job fails: its rows are
not stored, so the job is \field{failed}; a backfill job is run again from its
transcript, so its term is \emph{needs no spool}.

\begin{schema}{StoreFailsObservedCommittingCore}
\Delta Run \\ j? : JOB
\where
fixed = 1 \\
abl \neq 6 \\
((ck~j? = zfalse \land ts~j? = tpre) \lor (ck~j? = ztrue \land ts~j? = tcommit)) \\
loc~j? = out \\
st~j? = scommitting \\
tl~j? = tunconf \\
lc = zfalse \\
st' = st \oplus \{ j? \mapsto sfailed \} \\
loc' = loc \\
ts' = ts \oplus \{ j? \mapsto tstop \} \\
ab' = ab \\
sp' = sp \\
tl' = tl \oplus \{ j? \mapsto tnospool \} \\
rs' = rs \oplus \{ j? \mapsto rnospool \} \\
fl' = fl \\
fc' = fc \\
rp' = rp \\
cnt' = cnt \oplus \{ tunconf \mapsto cnt~tunconf - 1, tnospool \mapsto cnt~tnospool + 1 \} \\
ph' = ph \\
lc' = lc \\
bx' = bx
\end{schema}

\begin{zed}
StoreFailsObservedCommitting \defs StoreFailsObservedCommittingCore \land Latch
\end{zed}

\emph{Fixed build.} The same, before the abort has begun the job.

\begin{schema}{StoreFailsObservedUnbegunCore}
\Delta Run \\ j? : JOB
\where
fixed = 1 \\
((ck~j? = zfalse \land ts~j? = tpre) \lor (ck~j? = ztrue \land ts~j? = tcommit)) \\
loc~j? = out \\
st~j? = scommitting \\
tl~j? = tnone \\
lc = zfalse \\
st' = st \oplus \{ j? \mapsto sfailed \} \\
loc' = loc \\
ts' = ts \oplus \{ j? \mapsto tstop \} \\
ab' = ab \\
sp' = sp \\
tl' = tl \\
rs' = rs \oplus \{ j? \mapsto rterm \} \\
fl' = fl \\
fc' = fc \\
rp' = rp \\
cnt' = cnt \\
ph' = ph \\
lc' = lc \\
bx' = bx
\end{schema}

\begin{zed}
StoreFailsObservedUnbegun \defs StoreFailsObservedUnbegunCore \land Latch
\end{zed}

\emph{Fixed build.} A detached store that fails on a job whose status already
says it did not store (failed) or is the producer's: nothing to correct.

\begin{schema}{StoreFailsDetachedOtherCore}
\Delta Run \\ j? : JOB
\where
fixed = 1 \\
((ck~j? = zfalse \land ts~j? = tpre) \lor (ck~j? = ztrue \land ts~j? = tcommit)) \\
loc~j? = out \\
st~j? \neq scommitting \\
lc = zfalse \\
st' = st \\
loc' = loc \\
ts' = ts \oplus \{ j? \mapsto tstop \} \\
ab' = ab \\
sp' = sp \\
tl' = tl \\
rs' = rs \\
fl' = fl \\
fc' = fc \\
rp' = rp \\
cnt' = cnt \\
ph' = ph \\
lc' = lc \\
bx' = bx
\end{schema}

\begin{zed}
StoreFailsDetachedOther \defs StoreFailsDetachedOtherCore \land Latch
\end{zed}


% ===================================================================
\section{The worker's terminal}
\label{sec:terminal}
% ===================================================================

\field{TaskTerminal.recording} and \field{RefusalAwareTerminal} record how the
body ended. Before the abort the body's exit is \field{Finish}; the cancel
reaches the body through the \field{CancelInBody} family.

\emph{Old build.} The worker's terminal records \field{completed}
(\field{RefusalAwareTerminal.complete}) with no claim and over any status.

\begin{schema}{FinishStoredOldCore}
\Delta Run \\ j? : JOB
\where
fixed = 0 \\
ph = pre \\
loc~j? = body \\
ts~j? = tstored \\
st' = st \oplus \{ j? \mapsto scompleted \} \\
loc' = loc \oplus \{ j? \mapsto out \} \\
ts' = ts \\
ab' = ab \\
sp' = sp \\
tl' = tl \\
rs' = rs \\
fl' = fl \\
fc' = fc \\
rp' = rp \\
cnt' = cnt \\
ph' = ph \\
lc' = lc \\
bx' = bx
\end{schema}

\begin{zed}
FinishStoredOld \defs FinishStoredOldCore \land LatchClaim
\end{zed}

\emph{Old build.} A store that ended without rows (a refusal, the checkpoint's
\field{JobWithdrawnError}, an exception) records \field{failed} with no claim
and over any status, a producer's \field{aborted} included.

\begin{schema}{FinishStopOldCore}
\Delta Run \\ j? : JOB
\where
fixed = 0 \\
ph = pre \\
loc~j? = body \\
ts~j? = tstop \\
st' = st \oplus \{ j? \mapsto sfailed \} \\
loc' = loc \oplus \{ j? \mapsto out \} \\
ts' = ts \\
ab' = ab \\
sp' = sp \\
tl' = tl \\
rs' = rs \oplus \{ j? \mapsto rterm \} \\
fl' = fl \\
fc' = fc \\
rp' = rp \\
cnt' = cnt \\
ph' = ph \\
lc' = lc \\
bx' = bx
\end{schema}

\begin{zed}
FinishStopOld \defs FinishStopOldCore \land LatchClaim
\end{zed}

\emph{Fixed build.} The terminal records \field{completed} under the claim, and
only over a job the worker still owns (running or committing).

\begin{schema}{FinishStoredCore}
\Delta Run \\ j? : JOB
\where
fixed = 1 \\
ph = pre \\
loc~j? = body \\
ts~j? = tstored \\
st~j? \in \{ srunning, scommitting \} \\
st' = st \oplus \{ j? \mapsto scompleted \} \\
loc' = loc \oplus \{ j? \mapsto out \} \\
ts' = ts \\
ab' = ab \\
sp' = sp \\
tl' = tl \\
rs' = rs \\
fl' = fl \\
fc' = fc \\
rp' = rp \\
cnt' = cnt \\
ph' = ph \\
lc' = lc \\
bx' = bx
\end{schema}

\begin{zed}
FinishStored \defs FinishStoredCore \land Latch
\end{zed}

\emph{Fixed build.} A store without rows records \field{failed} under the claim,
over a running or committing job (the store is over, so the commit is settled).

\begin{schema}{FinishStopCore}
\Delta Run \\ j? : JOB
\where
fixed = 1 \\
ph = pre \\
loc~j? = body \\
ts~j? = tstop \\
st~j? \in \{ srunning, scommitting \} \\
st' = st \oplus \{ j? \mapsto sfailed \} \\
loc' = loc \oplus \{ j? \mapsto out \} \\
ts' = ts \\
ab' = ab \\
sp' = sp \\
tl' = tl \\
rs' = rs \oplus \{ j? \mapsto rterm \} \\
fl' = fl \\
fc' = fc \\
rp' = rp \\
cnt' = cnt \\
ph' = ph \\
lc' = lc \\
bx' = bx
\end{schema}

\begin{zed}
FinishStop \defs FinishStopCore \land Latch
\end{zed}

\emph{Fixed build.} The producer withdrew the job (\field{aborted}); the body's
exit leaves that status alone.

\begin{schema}{FinishKeptCore}
\Delta Run \\ j? : JOB
\where
fixed = 1 \\
ph = pre \\
loc~j? = body \\
ts~j? \in \{ tstored, tstop \} \\
st~j? = saborted \\
st' = st \\
loc' = loc \oplus \{ j? \mapsto out \} \\
ts' = ts \\
ab' = ab \\
sp' = sp \\
tl' = tl \\
rs' = rs \\
fl' = fl \\
fc' = fc \\
rp' = rp \\
cnt' = cnt \\
ph' = ph \\
lc' = lc \\
bx' = bx
\end{schema}

\begin{zed}
FinishKept \defs FinishKeptCore \land Latch
\end{zed}


% ===================================================================
\section{The abort: taking the jobs}
\label{sec:take}
% ===================================================================

\field{CollectionWorker.abort} cancels the worker task and takes the queued
jobs; the in-flight job reaches the abort through the worker's handler.

\field{CollectionWorker.abort}: the worker task is cancelled and every job still
waiting in its queue is taken, in one synchronous step.

\begin{schema}{AbortStartCore}
\Delta Run
\where
ph = pre \\
st' = st \\
loc' = \{ j : JOB; v : LOC | (loc~j = inq \land v = qtaken) \lor (loc~j \neq inq \land v = loc~j) @ j \mapsto v \} \\
ts' = ts \\
ab' = \{ j : JOB; v : AB | (loc~j = inq \land v = atake) \lor (loc~j \neq inq \land v = ab~j) @ j \mapsto v \} \\
sp' = sp \\
tl' = tl \\
rs' = rs \\
fl' = fl \\
fc' = fc \\
rp' = rp \\
cnt' = cnt \\
ph' = started \\
lc' = lc \\
bx' = bx
\end{schema}

\begin{zed}
AbortStart \defs AbortStartCore \land Latch
\end{zed}

The cancel reaches the in-flight job while it waits for the embed gate: its
terminal never ran, and the worker's handler takes the job.

\begin{schema}{CancelAtGateCore}
\Delta Run \\ j? : JOB
\where
ph = started \\
loc~j? = gate \\
st' = st \\
loc' = loc \oplus \{ j? \mapsto out \} \\
ts' = ts \\
ab' = ab \oplus \{ j? \mapsto atake \} \\
sp' = sp \\
tl' = tl \\
rs' = rs \\
fl' = fl \oplus \{ j? \mapsto ztrue \} \\
fc' = fc \\
rp' = rp \\
cnt' = cnt \\
ph' = ph \\
lc' = lc \\
bx' = bx
\end{schema}

\begin{zed}
CancelAtGate \defs CancelAtGateCore \land Latch
\end{zed}

\emph{Old build.} The cancel reaches the job inside its terminal
(\field{TaskTerminal.recording}), which records \field{failed} with no claim and
over whatever the status is: a committing job, a producer's \field{aborted}, and
a store that has in fact already landed.

\begin{schema}{CancelInBodyOldCore}
\Delta Run \\ j? : JOB
\where
fixed = 0 \\
ph = started \\
loc~j? = body \\
st' = st \oplus \{ j? \mapsto sfailed \} \\
loc' = loc \oplus \{ j? \mapsto out \} \\
ts' = ts \\
ab' = ab \oplus \{ j? \mapsto atake \} \\
sp' = sp \\
tl' = tl \\
rs' = rs \oplus \{ j? \mapsto rterm \} \\
fl' = fl \oplus \{ j? \mapsto ztrue \} \\
fc' = fc \\
rp' = rp \\
cnt' = cnt \\
ph' = ph \\
lc' = lc \\
bx' = bx
\end{schema}

\begin{zed}
CancelInBodyOld \defs CancelInBodyOldCore \land LatchClaim
\end{zed}

\emph{Fixed build.} The cancel reaches the job inside its terminal. A running
job whose store has not landed is \field{failed}, provisionally, under the
claim.

\begin{schema}{CancelInBodyFailsCore}
\Delta Run \\ j? : JOB
\where
fixed = 1 \\
ph = started \\
loc~j? = body \\
((st~j? = srunning \land ts~j? \neq tstored) \lor (st~j? = scommitting \land ts~j? = tstop)) \\
st' = st \oplus \{ j? \mapsto sfailed \} \\
loc' = loc \oplus \{ j? \mapsto out \} \\
ts' = ts \\
ab' = ab \oplus \{ j? \mapsto atake \} \\
sp' = sp \\
tl' = tl \\
rs' = rs \oplus \{ j? \mapsto rterm \} \\
fl' = fl \oplus \{ j? \mapsto ztrue \} \\
fc' = fc \\
rp' = rp \\
cnt' = cnt \\
ph' = ph \\
lc' = lc \\
bx' = bx
\end{schema}

\begin{zed}
CancelInBodyFails \defs CancelInBodyFailsCore \land Latch
\end{zed}

\emph{Fixed build.} The store had landed before the cancel and the cancelled
await discarded its result: the terminal reads the job's own future, sees the
rows, and records \field{completed}.

\begin{schema}{CancelInBodyStoredCore}
\Delta Run \\ j? : JOB
\where
fixed = 1 \\
ph = started \\
loc~j? = body \\
ts~j? = tstored \\
st~j? \in \{ srunning, scommitting \} \\
st' = st \oplus \{ j? \mapsto scompleted \} \\
loc' = loc \oplus \{ j? \mapsto out \} \\
ts' = ts \\
ab' = ab \oplus \{ j? \mapsto atake \} \\
sp' = sp \\
tl' = tl \\
rs' = rs \\
fl' = fl \oplus \{ j? \mapsto ztrue \} \\
fc' = fc \\
rp' = rp \\
cnt' = cnt \\
ph' = ph \\
lc' = lc \\
bx' = bx
\end{schema}

\begin{zed}
CancelInBodyStored \defs CancelInBodyStoredCore \land Latch
\end{zed}

\emph{Fixed build.} A committing job whose store is still running stays
\field{committing} (it may yet store), and a producer's \field{aborted} stays.

\begin{schema}{CancelInBodyKeptCore}
\Delta Run \\ j? : JOB
\where
fixed = 1 \\
ph = started \\
loc~j? = body \\
(st~j? = saborted \lor (st~j? = scommitting \land ts~j? = tcommit)) \\
st' = st \\
loc' = loc \oplus \{ j? \mapsto out \} \\
ts' = ts \\
ab' = ab \oplus \{ j? \mapsto atake \} \\
sp' = sp \\
tl' = tl \\
rs' = rs \\
fl' = fl \oplus \{ j? \mapsto ztrue \} \\
fc' = fc \\
rp' = rp \\
cnt' = cnt \\
ph' = ph \\
lc' = lc \\
bx' = bx
\end{schema}

\begin{zed}
CancelInBodyKept \defs CancelInBodyKeptCore \land Latch
\end{zed}


% ===================================================================
\section{The abort: beginning, deciding, recording}
\label{sec:record}
% ===================================================================

\field{JobAbort} fails each taken job with a provisional reason and counts it
pending (\field{Begin}), decides what it needs (\field{Decide}), and records the
outcome when its spool thread ends (\field{Record}), when the in-flight wait
expires (\field{Deadline}), or later (\field{Late}).

\emph{Old build.} \field{JobAbort._begin}: every taken job is failed with the
provisional reason, whatever \field{fail\_aborted} said, and counted pending. The
kept set lacks \field{aborted}, so a producer's \field{aborted} is overwritten.

\begin{schema}{BeginOldCore}
\Delta Run \\ j? : JOB
\where
fixed = 0 \\
ab~j? = atake \\
(st~j? \in \{ swithdrawn, scommitting, scompleted \} \implies st' = st) \land (st~j? \notin \{ swithdrawn, scommitting, scompleted \} \implies st' = st \oplus \{ j? \mapsto sfailed \}) \\
loc' = loc \\
ts' = ts \\
ab' = ab \oplus \{ j? \mapsto abegun \} \\
sp' = sp \\
tl' = tl \oplus \{ j? \mapsto tpend \} \\
(st~j? \in \{ swithdrawn, scommitting, scompleted \} \implies rs' = rs) \land (st~j? \notin \{ swithdrawn, scommitting, scompleted \} \implies rs' = rs \oplus \{ j? \mapsto rprov \}) \\
fl' = fl \\
fc' = fc \\
rp' = rp \\
cnt' = cnt \oplus \{ tnone \mapsto cnt~tnone - 1, tpend \mapsto cnt~tpend + 1 \} \\
ph' = ph \\
lc' = lc \\
bx' = bx
\end{schema}

\begin{zed}
BeginOld \defs BeginOldCore \land Latch
\end{zed}

\emph{Fixed build.} A taken job the abort must keep (withdrawn, aborted by its
producer, or stored) is counted \emph{kept} at once: it is neither failed nor
spooled.

\begin{schema}{BeginKeptCore}
\Delta Run \\ j? : JOB
\where
fixed = 1 \\
ab~j? = atake \\
st~j? \in \{ swithdrawn, saborted, scompleted \} \\
st' = st \\
loc' = loc \\
ts' = ts \\
ab' = ab \oplus \{ j? \mapsto arec \} \\
sp' = sp \\
tl' = tl \oplus \{ j? \mapsto tkept \} \\
rs' = rs \\
fl' = fl \\
fc' = fc \\
rp' = rp \\
cnt' = cnt \oplus \{ tnone \mapsto cnt~tnone - 1, tkept \mapsto cnt~tkept + 1 \} \\
ph' = ph \\
lc' = lc \\
bx' = bx
\end{schema}

\begin{zed}
BeginKept \defs BeginKeptCore \land Latch
\end{zed}

\emph{Fixed build.} A committing job may yet store; it keeps its status and is
counted \emph{unconfirmed} until its store thread ends.

\begin{schema}{BeginUnconfirmedCore}
\Delta Run \\ j? : JOB
\where
fixed = 1 \\
ab~j? = atake \\
st~j? = scommitting \\
st' = st \\
loc' = loc \\
ts' = ts \\
ab' = ab \oplus \{ j? \mapsto arec \} \\
sp' = sp \\
tl' = tl \oplus \{ j? \mapsto tunconf \} \\
rs' = rs \\
fl' = fl \\
fc' = fc \\
rp' = rp \\
cnt' = cnt \oplus \{ tnone \mapsto cnt~tnone - 1, tunconf \mapsto cnt~tunconf + 1 \} \\
ph' = ph \\
lc' = lc \\
bx' = bx
\end{schema}

\begin{zed}
BeginUnconfirmed \defs BeginUnconfirmedCore \land Latch
\end{zed}

\emph{Fixed build.} Any other taken job is failed with the provisional reason,
under the claim, and counted pending.

\begin{schema}{BeginFailCore}
\Delta Run \\ j? : JOB
\where
fixed = 1 \\
ab~j? = atake \\
st~j? \in \{ squeued, srunning, sfailed \} \\
\lnot (st~j? = sfailed \land ts~j? = tstop) \\
st' = st \oplus \{ j? \mapsto sfailed \} \\
loc' = loc \\
ts' = ts \\
ab' = ab \oplus \{ j? \mapsto abegun \} \\
sp' = sp \\
tl' = tl \oplus \{ j? \mapsto tpend \} \\
rs' = rs \oplus \{ j? \mapsto rprov \} \\
fl' = fl \\
fc' = fc \\
rp' = rp \\
cnt' = cnt \oplus \{ tnone \mapsto cnt~tnone - 1, tpend \mapsto cnt~tpend + 1 \} \\
ph' = ph \\
lc' = lc \\
bx' = bx
\end{schema}

\begin{zed}
BeginFail \defs BeginFailCore \land Latch
\end{zed}

\emph{Fixed build.} A taken job whose own store has already ended without rows (it
refused, failed, or stopped at the checkpoint) is \emph{settled}: it keeps the
store's reason and is counted \emph{needs no spool}. It is not spooled (its content
was refused or is rerunnable), and it is not counted \emph{kept}, because \emph{kept}
with a \field{failed} status would break I9.

\begin{schema}{BeginSettledCore}
\Delta Run \\ j? : JOB
\where
fixed = 1 \\
abl \neq 9 \\
ab~j? = atake \\
st~j? = sfailed \\
ts~j? = tstop \\
st' = st \\
loc' = loc \\
ts' = ts \\
ab' = ab \oplus \{ j? \mapsto arec \} \\
sp' = sp \\
tl' = tl \oplus \{ j? \mapsto tnospool \} \\
rs' = rs \\
fl' = fl \\
fc' = fc \\
rp' = rp \\
cnt' = cnt \oplus \{ tnone \mapsto cnt~tnone - 1, tnospool \mapsto cnt~tnospool + 1 \} \\
ph' = ph \\
lc' = lc \\
bx' = bx
\end{schema}

\begin{zed}
BeginSettled \defs BeginSettledCore \land Latch
\end{zed}

\emph{Seeded fault 9.} The settled job is counted \emph{kept}.

\begin{schema}{BeginSettledKeptCore}
\Delta Run \\ j? : JOB
\where
fixed = 1 \\
abl = 9 \\
ab~j? = atake \\
st~j? = sfailed \\
ts~j? = tstop \\
st' = st \\
loc' = loc \\
ts' = ts \\
ab' = ab \oplus \{ j? \mapsto arec \} \\
sp' = sp \\
tl' = tl \oplus \{ j? \mapsto tkept \} \\
rs' = rs \\
fl' = fl \\
fc' = fc \\
rp' = rp \\
cnt' = cnt \oplus \{ tnone \mapsto cnt~tnone - 1, tkept \mapsto cnt~tkept + 1 \} \\
ph' = ph \\
lc' = lc \\
bx' = bx
\end{schema}

\begin{zed}
BeginSettledKept \defs BeginSettledKeptCore \land Latch
\end{zed}

\emph{Seeded fault 4.} The fixed build's kept set lacks \field{aborted}.

\begin{schema}{BeginAbortedOverwrittenCore}
\Delta Run \\ j? : JOB
\where
fixed = 1 \\
abl = 4 \\
ab~j? = atake \\
st~j? = saborted \\
st' = st \oplus \{ j? \mapsto sfailed \} \\
loc' = loc \\
ts' = ts \\
ab' = ab \oplus \{ j? \mapsto abegun \} \\
sp' = sp \\
tl' = tl \oplus \{ j? \mapsto tpend \} \\
rs' = rs \oplus \{ j? \mapsto rprov \} \\
fl' = fl \\
fc' = fc \\
rp' = rp \\
cnt' = cnt \oplus \{ tnone \mapsto cnt~tnone - 1, tpend \mapsto cnt~tpend + 1 \} \\
ph' = ph \\
lc' = lc \\
bx' = bx
\end{schema}

\begin{zed}
BeginAbortedOverwritten \defs BeginAbortedOverwrittenCore \land Latch
\end{zed}

\field{AbortReason.of} for a job whose source outlives it (a backfill transcript):
no spool, its term is \emph{needs no spool}. The old build does this to a kept
job too.

\begin{schema}{DecideRerunCore}
\Delta Run \\ j? : JOB
\where
ab~j? = abegun \\
sk~j? = skrerun \\
(st~j? \in \{ swithdrawn, scommitting, scompleted \} \implies st' = st) \land (st~j? \notin \{ swithdrawn, scommitting, scompleted \} \implies st' = st \oplus \{ j? \mapsto sfailed \}) \\
loc' = loc \\
ts' = ts \\
ab' = ab \oplus \{ j? \mapsto arec \} \\
sp' = sp \\
tl' = tl \oplus \{ j? \mapsto tnospool \} \\
(st~j? \in \{ swithdrawn, scommitting, scompleted \} \implies rs' = rs) \land (st~j? \notin \{ swithdrawn, scommitting, scompleted \} \implies rs' = rs \oplus \{ j? \mapsto rnospool \}) \\
fl' = fl \\
fc' = fc \\
rp' = rp \\
cnt' = cnt \oplus \{ tpend \mapsto cnt~tpend - 1, tnospool \mapsto cnt~tnospool + 1 \} \\
ph' = ph \\
lc' = lc \\
bx' = bx
\end{schema}

\begin{zed}
DecideRerun \defs DecideRerunCore \land Latch
\end{zed}

A job that needs a spool record starts its write on a daemon thread. The old
build starts it for a kept job too.

\begin{schema}{DecideSpoolCore}
\Delta Run \\ j? : JOB
\where
ab~j? = abegun \\
sk~j? = skspool \\
st' = st \\
loc' = loc \\
ts' = ts \\
ab' = ab \oplus \{ j? \mapsto aspool \} \\
sp' = sp \oplus \{ j? \mapsto srun \} \\
tl' = tl \\
rs' = rs \\
fl' = fl \\
fc' = fc \\
rp' = rp \\
cnt' = cnt \\
ph' = ph \\
lc' = lc \\
bx' = bx
\end{schema}

\begin{zed}
DecideSpool \defs DecideSpoolCore \land Latch
\end{zed}

The spool thread ends: the file is on disk (\field{sok}), the write failed
(\field{sfail}), or the record could not be built (\field{sunb}). Whether the
abort is still waiting is a separate matter.

\begin{schema}{SpoolEndsCore}
\Delta Run \\ j? : JOB \\ o? : SP
\where
sp~j? = srun \\
o? \in \{ sok, sfail, sunb \} \\
(fixed = 0 \lor abl = 5 \lor lc = ztrue \lor st~j? \neq scompleted) \\
st' = st \\
loc' = loc \\
ts' = ts \\
ab' = ab \\
sp' = sp \oplus \{ j? \mapsto o? \} \\
tl' = tl \\
rs' = rs \\
fl' = fl \\
fc' = fc \\
rp' = rp \\
cnt' = cnt \\
ph' = ph \\
lc' = lc \\
bx' = bx
\end{schema}

\begin{zed}
SpoolEnds \defs SpoolEndsCore \land Latch
\end{zed}

The outcome reaches the loop within the wait (\field{JobAbort.\_record}): the
job's reason and its term are written together. A status the abort does not own
keeps its reason (\field{fail\_aborted} refuses); the old build counts it
anyway.

\begin{schema}{RecordCore}
\Delta Run \\ j? : JOB
\where
ab~j? = aspool \\
sp~j? \in \{ sok, sfail, sunb \} \\
fc~j? = zfalse \\
tl~j? = tpend \\
(st~j? \in \{ swithdrawn, scommitting, scompleted \} \implies st' = st) \land (st~j? \notin \{ swithdrawn, scommitting, scompleted \} \implies st' = st \oplus \{ j? \mapsto sfailed \}) \\
loc' = loc \\
ts' = ts \\
ab' = ab \oplus \{ j? \mapsto arec \} \\
sp' = sp \\
tl' = tl \oplus \{ j? \mapsto TermOf~(sp~j?) \} \\
(st~j? \in \{ swithdrawn, scommitting, scompleted \} \implies rs' = rs) \land (st~j? \notin \{ swithdrawn, scommitting, scompleted \} \implies rs' = rs \oplus \{ j? \mapsto RsOf~(sp~j?) \}) \\
fl' = fl \\
fc' = fc \\
rp' = rp \\
cnt' = cnt \oplus \{ tpend \mapsto cnt~tpend - 1, TermOf~(sp~j?) \mapsto cnt~(TermOf~(sp~j?)) + 1 \} \\
ph' = ph \\
lc' = lc \\
bx' = bx
\end{schema}

\begin{zed}
Record \defs RecordCore \land Latch
\end{zed}

\field{JobAbort.fail\_in\_flight}: the in-flight job's spool is still unfinished
at \field{SPOOL\_WAIT\_S}; the job is counted \emph{unfinished} and the late
callback is armed.

\begin{schema}{DeadlineCore}
\Delta Run \\ j? : JOB
\where
ab~j? = aspool \\
fl~j? = ztrue \\
tl~j? = tpend \\
st' = st \\
loc' = loc \\
ts' = ts \\
ab' = ab \oplus \{ j? \mapsto arec \} \\
sp' = sp \\
tl' = tl \oplus \{ j? \mapsto tunfin \} \\
rs' = rs \\
fl' = fl \\
fc' = fc \\
rp' = rp \\
cnt' = cnt \oplus \{ tpend \mapsto cnt~tpend - 1, tunfin \mapsto cnt~tunfin + 1 \} \\
ph' = ph \\
lc' = lc \\
bx' = bx
\end{schema}

\begin{zed}
Deadline \defs DeadlineCore \land Latch
\end{zed}

\field{JobAbort.\_ended\_late}: a spool the abort stopped waiting for ends. The
job's reason and its term move together, once.

\begin{schema}{LateCore}
\Delta Run \\ j? : JOB
\where
ab~j? = arec \\
tl~j? = tunfin \\
sp~j? \in \{ sok, sfail, sunb \} \\
fc~j? = zfalse \\
abl \neq 2 \\
lc = zfalse \\
(st~j? \in \{ swithdrawn, scommitting, scompleted \} \implies st' = st) \land (st~j? \notin \{ swithdrawn, scommitting, scompleted \} \implies st' = st \oplus \{ j? \mapsto sfailed \}) \\
loc' = loc \\
ts' = ts \\
ab' = ab \\
sp' = sp \\
tl' = tl \oplus \{ j? \mapsto TermOf~(sp~j?) \} \\
(st~j? \in \{ swithdrawn, scommitting, scompleted \} \implies rs' = rs) \land (st~j? \notin \{ swithdrawn, scommitting, scompleted \} \implies rs' = rs \oplus \{ j? \mapsto RsOf~(sp~j?) \}) \\
fl' = fl \\
fc' = fc \\
rp' = rp \\
cnt' = cnt \oplus \{ tunfin \mapsto cnt~tunfin - 1, TermOf~(sp~j?) \mapsto cnt~(TermOf~(sp~j?)) + 1 \} \\
ph' = ph \\
lc' = lc \\
bx' = bx
\end{schema}

\begin{zed}
Late \defs LateCore \land Latch
\end{zed}

\emph{Seeded fault 2.} The late correction moves the term and not the reason.

\begin{schema}{LateNoReasonCore}
\Delta Run \\ j? : JOB
\where
ab~j? = arec \\
tl~j? = tunfin \\
sp~j? \in \{ sok, sfail, sunb \} \\
fc~j? = zfalse \\
abl = 2 \\
lc = zfalse \\
st' = st \\
loc' = loc \\
ts' = ts \\
ab' = ab \\
sp' = sp \\
tl' = tl \oplus \{ j? \mapsto TermOf~(sp~j?) \} \\
rs' = rs \\
fl' = fl \\
fc' = fc \\
rp' = rp \\
cnt' = cnt \oplus \{ tunfin \mapsto cnt~tunfin - 1, TermOf~(sp~j?) \mapsto cnt~(TermOf~(sp~j?)) + 1 \} \\
ph' = ph \\
lc' = lc \\
bx' = bx
\end{schema}

\begin{zed}
LateNoReason \defs LateNoReasonCore \land Latch
\end{zed}

\emph{Old build.} The spool's future is cancelled while its thread still runs
(the loop is torn down): the wait reports \field{NOT\_WRITTEN} though the file may
yet land, and nothing corrects it.

\begin{schema}{SpoolCancelOldCore}
\Delta Run \\ j? : JOB
\where
fixed = 0 \\
ab~j? = aspool \\
fl~j? = ztrue \\
sp~j? = srun \\
fc~j? = zfalse \\
tl~j? = tpend \\
st' = st \\
loc' = loc \\
ts' = ts \\
ab' = ab \oplus \{ j? \mapsto arec \} \\
sp' = sp \\
tl' = tl \oplus \{ j? \mapsto tnotw \} \\
(st~j? \in \{ swithdrawn, scommitting, scompleted \} \implies rs' = rs) \land (st~j? \notin \{ swithdrawn, scommitting, scompleted \} \implies rs' = rs \oplus \{ j? \mapsto rnotw \}) \\
fl' = fl \\
fc' = fc \oplus \{ j? \mapsto ztrue \} \\
rp' = rp \\
cnt' = cnt \oplus \{ tpend \mapsto cnt~tpend - 1, tnotw \mapsto cnt~tnotw + 1 \} \\
ph' = ph \\
lc' = lc \\
bx' = bx
\end{schema}

\begin{zed}
SpoolCancelOld \defs SpoolCancelOldCore \land Latch
\end{zed}

\emph{Fixed build.} The same cancel is reported as \emph{unfinished}: the thread
still runs, so what reaches disk is unknown, and no later callback can see it.

\begin{schema}{SpoolCancelCore}
\Delta Run \\ j? : JOB
\where
fixed = 1 \\
ab~j? = aspool \\
fl~j? = ztrue \\
sp~j? = srun \\
fc~j? = zfalse \\
tl~j? = tpend \\
st' = st \\
loc' = loc \\
ts' = ts \\
ab' = ab \oplus \{ j? \mapsto arec \} \\
sp' = sp \\
tl' = tl \oplus \{ j? \mapsto tunfin \} \\
rs' = rs \\
fl' = fl \\
fc' = fc \oplus \{ j? \mapsto ztrue \} \\
rp' = rp \\
cnt' = cnt \oplus \{ tpend \mapsto cnt~tpend - 1, tunfin \mapsto cnt~tunfin + 1 \} \\
ph' = ph \\
lc' = lc \\
bx' = bx
\end{schema}

\begin{zed}
SpoolCancel \defs SpoolCancelCore \land Latch
\end{zed}

\emph{Old build.} The abort's gather is cancelled while a queued job's spool
thread runs. \field{\_one\_off\_loop} catches only \field{Exception}, so the
cancel escapes it, \field{\_record} never runs, and the job stays pending for good.

\begin{schema}{QueueCancelOldCore}
\Delta Run \\ j? : JOB
\where
fixed = 0 \\
ab~j? = aspool \\
fl~j? = zfalse \\
sp~j? = srun \\
fc~j? = zfalse \\
st' = st \\
loc' = loc \\
ts' = ts \\
ab' = ab \\
sp' = sp \\
tl' = tl \\
rs' = rs \\
fl' = fl \\
fc' = fc \oplus \{ j? \mapsto ztrue \} \\
rp' = rp \\
cnt' = cnt \\
ph' = ph \\
lc' = lc \\
bx' = bx
\end{schema}

\begin{zed}
QueueCancelOld \defs QueueCancelOldCore \land Latch
\end{zed}

\emph{Fixed build.} The same cancel records the job \emph{unfinished} and arms the
late callback. The spool thread's future is shielded, so the cancel does not
cancel it and \field{Late} can still correct the record.

\begin{schema}{QueueCancelCore}
\Delta Run \\ j? : JOB
\where
fixed = 1 \\
ab~j? = aspool \\
fl~j? = zfalse \\
sp~j? = srun \\
fc~j? = zfalse \\
tl~j? = tpend \\
lc = zfalse \\
st' = st \\
loc' = loc \\
ts' = ts \\
ab' = ab \oplus \{ j? \mapsto arec \} \\
sp' = sp \\
tl' = tl \oplus \{ j? \mapsto tunfin \} \\
rs' = rs \\
fl' = fl \\
fc' = fc \\
rp' = rp \\
cnt' = cnt \oplus \{ tpend \mapsto cnt~tpend - 1, tunfin \mapsto cnt~tunfin + 1 \} \\
ph' = ph \\
lc' = lc \\
bx' = bx
\end{schema}

\begin{zed}
QueueCancel \defs QueueCancelCore \land Latch
\end{zed}

Every job is settled: nothing is left to do. The self-loop keeps the settled
state from being a deadlock, so that a deadlock anywhere else is a finding.

\begin{schema}{QuiesceCore}
\Delta Run
\where
ph = started \\
\forall j : JOB @ (ab~j = arec \lor (ab~j = an \land loc~j = out)) \land ts~j \in \{ t0, tstored, tstop \} \\
st' = st \\
loc' = loc \\
ts' = ts \\
ab' = ab \\
sp' = sp \\
tl' = tl \\
rs' = rs \\
fl' = fl \\
fc' = fc \\
rp' = rp \\
cnt' = cnt \\
ph' = ph \\
lc' = lc \\
bx' = bx
\end{schema}

\begin{zed}
Quiesce \defs QuiesceCore \land Latch
\end{zed}


% ===================================================================
\section{The event loop closes}
\label{sec:close}
% ===================================================================

\field{IngestQueue.aclose} returns and \field{asyncio.run} closes the loop. A store
thread the abort detached may still be running: the observer that would have
settled it is cancelled with the loop. The fixed build drains first, bounded by
the shutdown budget, and reports what it could not wait for.

The shutdown budget for the drain runs out. Time is not modelled: this may fire
whenever the abort has started.

\begin{schema}{BudgetExpiresCore}
\Delta Run
\where
ph = started \\
bx = zfalse \\
st' = st \\
loc' = loc \\
ts' = ts \\
ab' = ab \\
sp' = sp \\
tl' = tl \\
rs' = rs \\
fl' = fl \\
fc' = fc \\
rp' = rp \\
cnt' = cnt \\
ph' = ph \\
lc' = lc \\
bx' = ztrue
\end{schema}

\begin{zed}
BudgetExpires \defs BudgetExpiresCore \land Latch
\end{zed}

\emph{Fixed build.} \field{IngestQueue.aclose} has settled the abort and waited
for the detached stores and the pending spool removals, at most until the budget
expired; then \field{asyncio.run} closes the loop. Whatever is still running is
reported (\field{rp}): a warning names its task, and the tally logged is the
corrected one. After this no observer runs.

\begin{schema}{LoopCloseCore}
\Delta Run
\where
fixed = 1 \\
abl \neq 8 \\
ph = started \\
lc = zfalse \\
\forall j : JOB @ ab~j = arec \lor (ab~j = an \land loc~j = out) \\
((\forall j : JOB @ ts~j \notin \{ tpre, tcommit \} \land \lnot (st~j = scompleted \land sp~j = srun)) \lor bx = ztrue) \\
st' = st \\
loc' = loc \\
ts' = ts \\
ab' = ab \\
sp' = sp \\
tl' = tl \\
rs' = rs \\
fl' = fl \\
fc' = fc \\
rp' = \{ j : JOB; v : ZBOOL | ((ts~j \in \{ tpre, tcommit \} \lor (st~j = scompleted \land sp~j = srun)) \land v = ztrue) \lor (\lnot (ts~j \in \{ tpre, tcommit \} \lor (st~j = scompleted \land sp~j = srun)) \land v = rp~j) @ j \mapsto v \} \\
cnt' = cnt \\
ph' = ph \\
lc' = ztrue \\
bx' = bx
\end{schema}

\begin{zed}
LoopClose \defs LoopCloseCore \land Latch
\end{zed}

\emph{Old build, and seeded fault 8.} The loop closes with no drain and no report:
a store still running is simply abandoned, and nobody says so.

\begin{schema}{LoopCloseOldCore}
\Delta Run
\where
ph = started \\
lc = zfalse \\
\forall j : JOB @ ab~j = arec \lor (ab~j = an \land loc~j = out) \\
(fixed = 0 \lor abl = 8) \\
st' = st \\
loc' = loc \\
ts' = ts \\
ab' = ab \\
sp' = sp \\
tl' = tl \\
rs' = rs \\
fl' = fl \\
fc' = fc \\
rp' = rp \\
cnt' = cnt \\
ph' = ph \\
lc' = ztrue \\
bx' = bx
\end{schema}

\begin{zed}
LoopCloseOld \defs LoopCloseOldCore \land Latch
\end{zed}

\emph{Fixed build.} The store ends after the loop has closed. No observer runs:
the status stays what the abort left, and the job was reported at the close.

\begin{schema}{StoreEndsUnobservedCore}
\Delta Run \\ j? : JOB
\where
fixed = 1 \\
lc = ztrue \\
((ck~j? = zfalse \land ts~j? = tpre) \lor (ck~j? = ztrue \land ts~j? = tcommit)) \\
loc~j? = out \\
st' = st \\
loc' = loc \\
ts' = ts \oplus \{ j? \mapsto tstored \} \\
ab' = ab \\
sp' = sp \\
tl' = tl \\
rs' = rs \\
fl' = fl \\
fc' = fc \\
rp' = rp \\
cnt' = cnt \\
ph' = ph \\
lc' = lc \\
bx' = bx
\end{schema}

\begin{zed}
StoreEndsUnobserved \defs StoreEndsUnobservedCore \land Latch
\end{zed}

\emph{Fixed build.} The same, for a store that ends without rows.

\begin{schema}{StoreFailsUnobservedCore}
\Delta Run \\ j? : JOB
\where
fixed = 1 \\
lc = ztrue \\
((ck~j? = zfalse \land ts~j? = tpre) \lor (ck~j? = ztrue \land ts~j? = tcommit)) \\
loc~j? = out \\
st' = st \\
loc' = loc \\
ts' = ts \oplus \{ j? \mapsto tstop \} \\
ab' = ab \\
sp' = sp \\
tl' = tl \\
rs' = rs \\
fl' = fl \\
fc' = fc \\
rp' = rp \\
cnt' = cnt \\
ph' = ph \\
lc' = lc \\
bx' = bx
\end{schema}

\begin{zed}
StoreFailsUnobserved \defs StoreFailsUnobservedCore \land Latch
\end{zed}


% ===================================================================
\section{Properties checked}
\label{sec:props}
% ===================================================================

For \field{fixed}~$=1$, \field{abl}~$=0$: ProB finds no reachable state with any
witness in \field{wit}, no deadlock, and from every state after the abort has
started a state in which every job is settled is reachable. A job is
\emph{settled} when it is out of the abort's hands (never taken and
\field{out}, or \field{arec}) and no store thread is running
(\field{AG(started} $\Rightarrow$ \field{EF(settled))}). For \field{fixed}~$=0$ and
for each seeded fault, the witness aimed at it is reachable; the report records the
shortest path ProB found for each.

\end{document}
